\chapter{Introdução}

%% Contextualização do tema geral
%% Demonstrar sua relevância e importância


A busca por alternativas sustentáveis para a substituição dos combustíveis fósseis estabeleceu a urgência de uma nova transição energética global, cenário no qual os materiais bidimensionais (2D) despontam como a fronteira da inovação tecnológica. Caracterizados por sua espessura de escala atômica e por intensos efeitos de confinamento quântico, os sistemas 2D exibem propriedades físico-químicas singulares, tais como elevada área superficial específica, alta mobilidade de portadores de carga e facilidade de funcionalização superficial. Essas características conferem aos materiais bidimensionais um papel protagonista no desenvolvimento de tecnologias energéticas de próxima geração, incluindo eletrocatalisadores eficientes para a produção de hidrogênio verde, eletrodos de alta densidade para baterias e supercapacitores, além de camadas ativas em dispositivos fotovoltaicos avançados.

Para responder à velocidade exigida por essa transição energética, é indispensável explorar de forma rápida e abrangente um vasto espaço químico e estrutural de combinações atômicas. Como a síntese e a caracterização em laboratório são processos morosos e de alto custo, a ciência computacional de materiais passou a recorrer a bancos de dados de alta vazão (\textit{high-throughput}), como o Computational 2D Materials Database (C2DB)~\cite{haastrup2018computational}. Nesse contexto, a aplicação de modelos de inteligência artificial e aprendizado de máquina emergiu como uma abordagem fundamental para acelerar o mapeamento de propriedades físico-químicas determinantes, como o bandgap, a estabilidade estrutural e a energia de formação, viabilizando o projeto racional de novos compostos bidimensionais antes mesmo da etapa experimental~\cite{van2026survey}.

Apesar das promessas das abordagens computacionais, a determinação precisa das propriedades eletrônicas e estruturais em sistemas periódicos esbarra na limitação imposta pelo severo custo computacional dos métodos ab initio. A simulação teórica primária dessas estruturas apoia-se tradicionalmente na Teoria do Funcional da Densidade (DFT), implementada em pacotes computacionais consagrados como o VASP (Vienna Ab initio Simulation Package)~\cite{kresse1996efficient}. Por envolverem a resolução iterativa das equações de Kohn-Sham para elétrons interagentes sob condições periódicas de contorno, esses cálculos apresentam uma complexidade computacional que escala de forma não linear com o número de átomos na célula unitária. Consequentemente, a dependência exclusiva de cálculos ab initio impõe um gargalo proibitivo de tempo e poder de processamento em supercomputadores, tornando inviável a varredura exaustiva de grandes bibliotecas de materiais~\cite{langer2022representations}.

Embora a substituição dos cálculos ab initio por modelos de aprendizado profundo ofereça predições ordens de magnitude mais rápidas, a escassez de dados rotulados de alta fidelidade compromete o desempenho e a generalização de redes neurais puramente supervisionadas. É justamente para superar essa limitação que se justifica a adoção do aprendizado contrastivo. Por meio da combinação de abordagens auto-supervisionadas e supervisionadas, o aprendizado contrastivo mapeia a estrutura do material em um espaço latente no qual amostras estruturalmente ou funcionalmente semelhantes são aproximadas, enquanto amostras distintas são afastadas~\cite{koker2022graph,moro2025multimodal}. Essa estratégia permite que o modelo extraia representações físicas robustas a partir de grandes volumes de dados não rotulados ou parcialmente rotulados, contornando a necessidade de gerar um número massivo de novos rótulos via DFT e aumentando a acurácia da predição de propriedades em materiais bidimensionais~\cite{van2026survey}.

Diante desse cenário, este trabalho investiga a aplicação de técnicas de aprendizado contrastivo para a extração de representações latentes de materiais bidimensionais e a subsequente predição de suas propriedades eletrônicas e termodinâmicas, como o \textit{band gap} e a energia de formação. A proposta busca avaliar como o alinhamento métrico no espaço latente contribui para melhorar a capacidade de generalização dos modelos preditivos em bases de dados computacionais de sistemas 2D, comparando o desempenho dessas representações com abordagens supervisionadas tradicionais.

